\chapter{Validation, Safety, and Reporting Checklists}
\label{app:validation}

\section{Plasma and magnetic validation}

\begin{enumerate}[label=$\square$]
\item Coil current, phase, and frequency are independently recorded.
\item Vacuum field mapping agrees with the electromagnetic model.
\item Plasma response is measured or modeled in three dimensions.
\item Density, temperature, and impurity changes are reported with the radiation result.
\item At least two diagnostics constrain the energetic-electron population.
\item Particle and energy balances close within stated uncertainty.
\item Collector and wall heat loads are measured or bounded.
\end{enumerate}

\section{Radiation-source validation}

\begin{enumerate}[label=$\square$]
\item Detector calibration covers the relevant energy and pulse range.
\item Pileup, dead time, saturation, and electromagnetic interference are tested.
\item Spectral unfolding includes response matrices and covariance.
\item Angular coverage is sufficient to identify redistribution.
\item Electron, neutron, and activation contributions are measured or bounded.
\item Baseline measurements are repeated before and after perturbation scans.
\end{enumerate}

\section{Dosimetry validation}

\begin{enumerate}[label=$\square$]
\item A traceable absorbed-dose method is used.
\item Detector energy dependence is addressed.
\item Pulse-dose-rate response is characterized.
\item Phantom geometry matches the transport model.
\item Depth-dose and lateral profiles are measured.
\item Monte Carlo calculations reproduce benchmark measurements.
\item Uncertainty is propagated from source to dose.
\end{enumerate}

\section{Biological validation}

\begin{enumerate}[label=$\square$]
\item Cell lines are authenticated and mycoplasma free.
\item Passage range, medium, oxygen, and plating efficiency are documented.
\item Exposure order is randomized and analysis blinded where feasible.
\item Independent biological replicates are used.
\item A conventional reference beam is included.
\item Both source-normalized and dose-matched comparisons are performed.
\item Physical dose is measured at the actual sample position.
\item Exclusion criteria and missing data are reported.
\end{enumerate}

\section{Radiation safety}

\begin{enumerate}[label=$\square$]
\item Work is authorized under the institutional radiation-safety program.
\item Shielding calculations and surveys cover normal and credible fault states.
\item Interlocks, warning indicators, and emergency shutdowns are tested.
\item Access control and operating procedures are documented.
\item Personnel dosimetry is used where required.
\item Activation and post-run dose rates are evaluated.
\item Magnetic-field, high-voltage, laser, vacuum, and biological hazards are integrated into one hazard analysis.
\end{enumerate}

\section{Minimum reporting table}

\begin{longtable}{p{0.26\textwidth}p{0.64\textwidth}}
\toprule
Category & Minimum report \\
\midrule
Plasma & density, $T_e$, $T_i$, composition, stored energy, duration \\
Magnetic control & equilibrium field, $\delta B/B$, mode, frequency, phase, geometry \\
Electron population & diagnostic, energy range, anisotropy, inferred distribution \\
Photon source & spectrum, angle, pulse structure, absolute normalization \\
Other radiation & electrons, ions, neutrons, activation \\
Transport & geometry version, material composition, code and physics data \\
Dosimetry & detector type, calibration, position, uncertainty, dose rate \\
Biology & model, endpoint, timing, replicates, dose matching, statistics \\
Reproducibility & data identifiers, scripts, software versions, quality flags \\
\bottomrule
\end{longtable}
